\chapter{Blockchains pour les traders}\label{ch:m3:blockchains-for-traders}

Un trader voit un \omterm{def:m3:blockchains-for-traders:stable}{stablecoin} coté avec une décote sur une plateforme décentralisée et une prime sur une
bourse. Il retire des \omterm{def:m3:blockchains-for-traders:key}{jetons} de la bourse pour capter l'écart. Le retrait attend dans une
file d'attente, est diffusé, attend un bloc, puis les confirmations que la plateforme réceptrice exige~;
les frais de la transaction, enchéris dans une enchère qu'elle ne voit pas, doublent pendant l'attente. Quand les \omterm{def:m3:blockchains-for-traders:key}{jetons}
arrivent, l'écart s'est refermé. Les marchés crypto se règlent sur des \omterm{def:m3:blockchains-for-traders:chain}{blockchains}, et chaque propriété d'une chaîne
qu'un ingénieur pourrait juger abstraite (comment les blocs sont produits, quand ils deviennent finals, ce qu'une transaction
coûte) est, pour un trader, une latence, un coût ou un risque. Ce chapitre expose ce dont un trader a besoin de
la mécanique~: registres et consensus, \omterm{def:m3:blockchains-for-traders:finality}{finalité}, transactions et leurs frais, le \omterm{def:m3:blockchains-for-traders:mempool}{mempool}, les
chaînes superposées bâties par-dessus, les \omterm{def:m3:blockchains-for-traders:bridge}{bridges} entre elles, et les \omterm{def:m3:blockchains-for-traders:stable}{stablecoins}, les dollars de ce marché.

\section{Registres, consensus et finalité}

\begin{definition}[Blockchain, protocole de consensus]\label{def:m3:blockchains-for-traders:chain}
Une \emph{blockchain}\index{blockchain} est un registre tenu par de nombreux ordinateurs indépendants sous la forme d'une chaîne de
blocs, chacun contenant un lot de transactions et une empreinte cryptographique du bloc précédent, de sorte que
modifier un bloc passé modifierait chacun des suivants. Un
\emph{protocole de consensus}\index{protocole de consensus} est la règle par laquelle ces ordinateurs s'accordent sur le bloc qui prolonge la chaîne.
\end{definition}

\begin{definition}[Preuve de travail, preuve d'enjeu, validateur]\label{def:m3:blockchains-for-traders:pow}
Avec la \emph{preuve de travail}\index{preuve de travail} (proof of work), le droit de proposer un bloc s'obtient en résolvant un
casse-tête de calcul coûteux, et la chaîne au travail cumulé le plus élevé est celle qui fait foi. En revanche, avec la
\emph{preuve d'enjeu}\index{preuve d'enjeu} (proof of stake), il est attribué parmi les \emph{validateurs}\index{validateur},
des participants qui verrouillent (mettent en jeu) le \omterm{def:m3:blockchains-for-traders:key}{jeton} natif de la chaîne à titre de collatéral, lequel peut être détruit s'ils
se comportent mal, et qui proposent et attestent des blocs.
\end{definition}

\begin{definition}[Finalité]\label{def:m3:blockchains-for-traders:finality}
La \emph{finalité}\index{finalité} est le point au-delà duquel une transaction ne peut plus être remise en cause~:
probabiliste sous \omterm{def:m3:blockchains-for-traders:pow}{preuve de travail}, où chaque bloc supplémentaire rend une remise en cause exponentiellement moins probable~;
économique sous les protocoles de \omterm{def:m3:blockchains-for-traders:pow}{preuve d'enjeu} qui finalisent des points de contrôle, où en remettre un en cause détruirait une
part importante de l'enjeu.
\end{definition}

Pour un trader, la \omterm{def:m3:blockchains-for-traders:finality}{finalité} s'apparente à un cycle de règlement-livraison (livre~1, chapitre~5) de forme différente. Une
plateforme ne crédite un dépôt qu'après un nombre de confirmations qu'elle choisit, arbitrant entre son propre risque de
dépôt annulé et l'attente de ses clients.

\begin{dated}{2026-09}{Temps entre blocs et finalité}\label{dat:m3:blockchains-for-traders:times}
Bitcoin vise un bloc toutes les dix minutes~; une convention répandue considère six confirmations, soit environ une heure,
comme final. Ethereum est passé de la \omterm{def:m3:blockchains-for-traders:pow}{preuve de travail} à la \omterm{def:m3:blockchains-for-traders:pow}{preuve d'enjeu} le 15 septembre 2022 (le Merge)~; le temps est
découpé en créneaux de 12 secondes et en époques de 32 créneaux, et un bloc est finalisé après deux époques, soit environ 12.8
minutes. Le bloc 26\,048\,449 d'Ethereum, produit à 16:17 UTC le 24 septembre 2026, avait une limite de 60 millions de \omterm{def:m3:blockchains-for-traders:gas}{gas},
dont il a utilisé 22.1~millions, à un \omterm{def:m3:blockchains-for-traders:gas}{frais de base} de 1.49~gwei.
\end{dated}

\section{Transactions, gas et marché des frais}

\begin{definition}[Clé privée, jeton]\label{def:m3:blockchains-for-traders:key}
Une \emph{clé privée}\index{clé privée} est le nombre secret dont la signature autorise les transactions depuis
une adresse~; quiconque la détient contrôle les actifs de cette adresse. Un \emph{jeton}\index{jeton} est une
unité transférable enregistrée par une \omterm{def:m3:blockchains-for-traders:chain}{blockchain}~: la monnaie native de la chaîne, ou une unité dont les soldes sont tenus par un programme
de la chaîne qui les tient à jour.
\end{definition}

\begin{definition}[Smart contract et oracle]\label{def:m3:blockchains-for-traders:contract}
Un \emph{smart contract}\index{smart contract} est un programme déployé sur une \omterm{def:m3:blockchains-for-traders:chain}{blockchain}, dont le code et l'état
sont publics et dont les fonctions peuvent être appelées par quiconque en envoyant une transaction~; la chaîne l'exécute et enregistre
le résultat. Un \emph{oracle}\index{oracle} est un mécanisme qui fait entrer des données venues de l'extérieur d'une chaîne (un cours,
l'issue d'un événement) dans un smart contract.
\end{definition}

\begin{definition}[Gas, frais de base, frais de priorité]\label{def:m3:blockchains-for-traders:gas}
Le \emph{gas}\index{gas} est l'unité dans laquelle Ethereum mesure le calcul et le stockage qu'une transaction
consomme~; chaque bloc a une limite de gas. Une transaction paie, par unité de gas, les \emph{frais de
base}\index{frais de base}, qui sont brûlés et fixés par le protocole, plus des \emph{frais de priorité}\index{frais de priorité}
(un pourboire) qu'elle choisit, et qui revient au proposant du bloc.
\end{definition}

\begin{proposition}[La règle des frais de base]\label{prop:m3:blockchains-for-traders:basefee}
Sous l'EIP-1559, la cible de \omterm{def:m3:blockchains-for-traders:gas}{gas} d'un bloc est la moitié de sa limite, et les \omterm{def:m3:blockchains-for-traders:gas}{frais de base} du bloc suivant sont ceux du bloc parent
multipliés par $1 + \tfrac18\,(g - g^{*})/g^{*}$, où $g$ est le \omterm{def:m3:blockchains-for-traders:gas}{gas} utilisé par le bloc parent et $g^{*}$ sa cible~:
au plus $+12.5\%$ après un bloc plein et $-12.5\%$ après un bloc vide. Après $k$ blocs pleins, les \omterm{def:m3:blockchains-for-traders:gas}{frais de base} valent
$1.125^k$ fois leur niveau de départ.
\end{proposition}

\begin{proof}
Il s'agit de la mise à jour de l'EIP avec $g = 2g^{*}$ (plein) ou $g = 0$ (vide)~; on itère sur $k$ blocs. L'arithmétique entière
du protocole arrondit chaque étape vers le bas (et fait qu'une hausse vaut au moins un wei).
\end{proof}

Vingt blocs pleins, soit quatre minutes sur Ethereum, multiplient les \omterm{def:m3:blockchains-for-traders:gas}{frais de base} par 10.5~: de 10~gwei à 105.45~gwei. Un
swap utilisant 150\,000 \omterm{def:m3:blockchains-for-traders:gas}{gas} avec un pourboire de 2~gwei coûte alors 48.35~\$ pour un cours de l'ether de 3\,000~\$ (illustratif),
contre 5.40~\$ auparavant. Les frais sont l'unité dans laquelle se mesure la latence d'un trader quand les blocs sont saturés.

\begin{omfigure}
\begin{tikzpicture}
\begin{axis}[width=11cm,height=4.8cm,axis y line*=left,
  xlabel={bloc},ylabel={frais de base, gwei},
  tick label style={font=\small},label style={font=\small},
  xmin=0,xmax=59,ymin=0,ymax=45,grid=major,grid style={gray!25},
  table/col sep=comma]
\addplot[omThm,very thick] table[x=block,y=base]{figdata/markets-3/14-blockchains-for-traders/shock.csv};\label{plot:m3:basefee}
\end{axis}
\begin{axis}[width=11cm,height=4.8cm,axis y line*=right,axis x line=none,
  ylabel={gas utilisé, \% de la limite},ylabel near ticks,
  tick label style={font=\small},label style={font=\small},
  xmin=0,xmax=59,ymin=0,ymax=110,
  legend columns=2,legend style={font=\footnotesize,at={(0.5,-0.3)},anchor=north,draw=none,fill=none,/tikz/every even column/.append style={column sep=8pt}},
  table/col sep=comma]
\addlegendimage{/pgfplots/refstyle=plot:m3:basefee}\addlegendentry{frais de base (gauche)}
\addplot[omDef,thick,dashed,const plot] table[x=block,y=share]{figdata/markets-3/14-blockchains-for-traders/shock.csv};\addlegendentry{gas utilisé (droite)}
\end{axis}
\end{tikzpicture}

\omcaption{Un choc de demande sous l'EIP-1559~: du bloc 5 au bloc 34, la demande de \omterm{def:m3:blockchains-for-traders:gas}{gas} est quatre fois plus élevée. Les blocs
se remplissent, les \omterm{def:m3:blockchains-for-traders:gas}{frais de base} augmentent de 12.5~\% par bloc jusqu'à ce que la demande retombe à la cible (la moitié de la limite), et décroissent
lorsque le choc prend fin. Demande illustrative~; arithmétique exacte du protocole. Données~: le tutoriel du chapitre.}\label{fig:m3:blockchains-for-traders:shock}
\end{omfigure}

\section{Le mempool}

\begin{definition}[Mempool]\label{def:m3:blockchains-for-traders:mempool}
Le \emph{mempool}\index{mempool} est l'ensemble des transactions signées que les nœuds ont reçues et relayées
mais qui ne figurent pas encore dans un bloc~: publiques, dans la plupart des chaînes, pour quiconque fait tourner un nœud.
\end{definition}

Une transaction présente dans le \omterm{def:m3:blockchains-for-traders:mempool}{mempool} public est un ordre que chacun peut voir avant son exécution. Cette visibilité
est à la racine de l'extraction présentée au \cref{ch:m3:maximal-extractable-value}~: un searcher qui voit un gros swap
en attente peut placer ses propres transactions avant et après elle. Cela signifie aussi que le transfert d'un trader est visible de ses
concurrents, et que son inclusion dépend de ce que les autres enchérissent~: les \omterm{def:m3:blockchains-for-traders:gas}{frais de priorité} sont une enchère pour la position,
et une transaction dont le pourboire est trop faible peut attendre plusieurs blocs.

\begin{omfigure}
\centering
\begin{tikzpicture}[font=\small,
  st/.style={draw,thick,rounded corners=2pt,align=center,minimum height=1cm,text width=2.6cm},
  arr/.style={-{Stealth[length=2.2mm]},thick}]
\node[st,draw=omDef,fill=omDef!8] (w) at (0,0) {le portefeuille signe\\avec la clé privée};
\node[st,draw=omProp,fill=omProp!8] (m) at (3.4,0) {mempool public\\(visible de tous)};
\node[st,draw=omThm,fill=omThm!8] (b) at (6.8,0) {incluse dans un\\bloc};
\node[st,draw=omExo,fill=omExo!8] (f) at (10.2,0) {confirmations,\\puis finalité};
\draw[arr] (w) -- (m); \draw[arr] (m) -- (b); \draw[arr] (b) -- (f);
\node[font=\footnotesize,anchor=north,align=center] at (5.1,-0.8) {secondes à minutes~:\\pourboires et congestion};
\node[font=\footnotesize,anchor=north,align=center] at (10.2,-0.8) {minutes à une heure~:\\le protocole};
\end{tikzpicture}

\omcaption{Le cycle de vie d'une transaction et l'endroit où passe le temps d'un trader~: l'attente de l'inclusion dépend des frais
payés par rapport à ceux des autres, l'attente de la \omterm{def:m3:blockchains-for-traders:finality}{finalité} dépend de la chaîne. Schéma de principe.}\label{fig:m3:blockchains-for-traders:life}
\end{omfigure}

\section{Couches de base, rollups, séquenceurs et bridges}

Les couches de base sont lentes et coûteuses, parce que chaque nœud exécute chaque transaction. L'essentiel de l'activité s'est
déplacé vers des chaînes construites au-dessus d'elles.

\begin{definition}[Rollup et séquenceur]\label{def:m3:blockchains-for-traders:rollup}
Un \emph{rollup}\index{rollup} est une chaîne qui exécute les transactions hors de la couche de base et les publie, en
lots compressés, sur la couche de base, qui garantit leurs données et, par des preuves de fraude ou de validité,
leur exactitude. Son \emph{séquenceur}\index{séquenceur} est l'opérateur qui ordonne les transactions du rollup
et produit ses blocs~; sur la plupart des rollups, il s'agit d'un opérateur unique.
\end{definition}

\begin{omfigure}
\centering
\begin{tikzpicture}[font=\small,
  l/.style={draw,thick,rounded corners=2pt,align=center,minimum height=0.9cm,text width=3.2cm},
  arr/.style={-{Stealth[length=2.2mm]},thick}]
\node[l,draw=omThm,fill=omThm!8] (seq) at (0,1.6) {séquenceur du rollup~:\\ordonne les transactions};
\node[l,draw=omDef,fill=omDef!8] (l1) at (0,-0.4) {couche de base~: stocke les lots,\\vérifie les preuves};
\node[l,draw=omProp,fill=omProp!8] (other) at (7.2,1.6) {une autre chaîne};
\node[l,draw=omExo,fill=omExo!8] (br) at (7.2,-0.4) {contrats de bridge~:\\verrouille ici, émet là-bas};
\draw[arr] (seq) -- node[right,font=\footnotesize] {lots compressés} (l1);
\draw[arr] (l1) -- node[above,font=\footnotesize] {verrouillage de jetons} (br);
\draw[arr] (br) -- node[right,font=\footnotesize] {émission d'une créance} (other);
\node[font=\footnotesize,anchor=east,align=right] at (-1.8,1.6) {les utilisateurs envoient\\des transactions};
\end{tikzpicture}

\omcaption{Un \omterm{def:m3:blockchains-for-traders:rollup}{rollup} ordonne les transactions hors de la couche de base et les publie là-bas par lots~; un \omterm{def:m3:blockchains-for-traders:bridge}{bridge} verrouille des \omterm{def:m3:blockchains-for-traders:key}{jetons}
sur une chaîne et émet sur une autre une créance sur ces \omterm{def:m3:blockchains-for-traders:key}{jetons}. Le \omterm{def:m3:blockchains-for-traders:rollup}{séquenceur} et les contrats ou les signataires du \omterm{def:m3:blockchains-for-traders:bridge}{bridge} sont les
points de confiance. Schéma de principe.}\label{fig:m3:blockchains-for-traders:rollup}
\end{omfigure}

Un \omterm{def:m3:blockchains-for-traders:rollup}{séquenceur} unique décide de l'ordre des transactions de son \omterm{def:m3:blockchains-for-traders:rollup}{rollup}, soit le même pouvoir, et la même
tentation, que celui d'un constructeur de blocs sur la couche de base~; il voit aussi les transactions avant qu'elles ne soient ordonnées. Dans un
\omterm{def:m3:blockchains-for-traders:rollup}{rollup}, les transactions sont peu coûteuses et rapides à inclure, mais elles ne sont définitives sur la couche de base que lorsque le lot l'est.

\begin{dated}{2026-09}{Qui séquence les rollups}\label{dat:m3:blockchains-for-traders:l2beat}
Sur les 101 projets de couche 2 recensés dans la synthèse de L2BEAT au 24 septembre 2026, 3 disposaient d'un ensemble de \omterm{def:m3:blockchains-for-traders:rollup}{séquenceurs} décentralisé~; pour 67, un utilisateur
pouvait séquencer lui-même ses transactions via la couche de base si le \omterm{def:m3:blockchains-for-traders:rollup}{séquenceur} tombait en panne, et pour 21 il n'existait aucun mécanisme de ce type.
\end{dated}

\begin{definition}[Bridge]\label{def:m3:blockchains-for-traders:bridge}
Un \emph{bridge}\index{bridge} transfère des \omterm{def:m3:blockchains-for-traders:key}{jetons} entre chaînes~: il les verrouille ou les brûle sur une chaîne et libère
ou émet un \omterm{def:m3:blockchains-for-traders:key}{jeton} correspondant sur l'autre, sur la foi d'une preuve ou d'un ensemble de signataires qui attestent
du verrouillage.
\end{definition}

Un \omterm{def:m3:blockchains-for-traders:bridge}{bridge} concentre les actifs dans un seul ensemble de contrats et de clés, ce qui en fait la plus grande cible de la crypto.
En mars 2022, des attaquants ayant obtenu la majorité des clés de signature du \omterm{def:m3:blockchains-for-traders:bridge}{bridge} Ronin ont retiré environ
600~millions de dollars de \omterm{def:m3:blockchains-for-traders:key}{jetons}~; en avril, l'US Treasury a attribué le vol à un groupe nord-coréen.

\section{Stablecoins~: émission et rachat}

\begin{definition}[Stablecoin, depeg]\label{def:m3:blockchains-for-traders:stable}
Un \emph{stablecoin}\index{stablecoin} est un \omterm{def:m3:blockchains-for-traders:key}{jeton} conçu pour maintenir une valeur fixe par rapport à une monnaie, généralement
le dollar, le plus souvent grâce à un émetteur qui détient des réserves (espèces, bons du Trésor, repo) et émet des \omterm{def:m3:blockchains-for-traders:key}{jetons} contre des dépôts
et les brûle contre des rachats. Un \emph{décrochage}\index{décrochage} (depeg) est une période pendant laquelle le \omterm{def:m3:blockchains-for-traders:key}{jeton} se négocie
à l'écart de sa valeur cible sur le marché secondaire.
\end{definition}

L'arrimage est tenu par l'arbitrage entre deux marchés. Sur le marché primaire, quelques institutions déposent des dollars
auprès de l'émetteur et reçoivent des \omterm{def:m3:blockchains-for-traders:key}{jetons}, ou restituent des \omterm{def:m3:blockchains-for-traders:key}{jetons} et reçoivent des dollars, à raison d'un pour un. Sur le
marché secondaire, tous les autres acteurs échangent des \omterm{def:m3:blockchains-for-traders:key}{jetons} sur les bourses. Lorsque le \omterm{def:m3:blockchains-for-traders:key}{jeton} se négocie sous le dollar, une institution
disposant d'un accès au marché primaire l'achète puis le fait racheter~; lorsqu'il se négocie au-dessus, elle émet et vend. L'arrimage ne vaut que ce que
vaut cet arbitrage, qui repose sur la valeur d'un dollar des réserves et sur l'ouverture des rachats.

\begin{omfigure}
\centering
\begin{tikzpicture}[font=\small,
  p/.style={draw,thick,rounded corners=2pt,align=center,minimum height=1cm,text width=2.5cm},
  arr/.style={-{Stealth[length=2.2mm]},thick}]
\node[p,draw=omDef,fill=omDef!8] (iss) at (0,0) {émetteur\\(réserves)};
\node[p,draw=omThm,fill=omThm!8] (inst) at (5.2,0) {institution disposant\\d'un accès primaire};
\node[p,draw=omProp,fill=omProp!8] (sec) at (10.4,0) {marché secondaire\\(bourses, pools)};
\draw[arr] ([yshift=5pt]inst.west) -- node[above,font=\footnotesize] {\$ entrant, émission} ([yshift=5pt]iss.east);
\draw[arr] ([yshift=-5pt]iss.east) -- node[below,font=\footnotesize] {destruction, \$ sortant} ([yshift=-5pt]inst.west);
\draw[arr,{Stealth[length=2.2mm]}-{Stealth[length=2.2mm]}] (inst) -- node[above,font=\footnotesize] {achète $<\$1$} node[below,font=\footnotesize] {vend $>\$1$} (sec);
\end{tikzpicture}

\omcaption{Comment un \omterm{def:m3:blockchains-for-traders:stable}{stablecoin} adossé à des réserves tient son ancrage~: les institutions disposant d'un accès primaire émettent et rachètent
au taux de un pour un et arbitrent l'écart sur le marché secondaire. Si le rachat s'interrompt ou si les réserves sont
mises en doute, rien ne maintient l'ancrage. Schéma.}\label{fig:m3:blockchains-for-traders:mint}
\end{omfigure}

\begin{dated}{2026-09}{Stablecoins}\label{dat:m3:blockchains-for-traders:stable}
L'offre totale de \omterm{def:m3:blockchains-for-traders:stable}{stablecoins} s'élevait à environ 303~milliards de dollars le 10 septembre 2026, dont environ 183~milliards pour l'USDT et
environ 74~milliards pour l'USDC (chiffres des agrégateurs). En mars 2023, Circle n'a pas pu retirer 3.3~milliards de dollars de réserves de l'USDC,
soit environ 8~\% d'entre elles, détenues auprès de la Silicon Valley Bank défaillante~; l'USDC est descendu jusqu'à environ 86~cents avant que
les réserves soient récupérées et l'ancrage rétabli. Les chercheurs de la Réserve fédérale relèvent que la plupart des détenteurs
ne peuvent pas racheter auprès de l'émetteur et ne peuvent que négocier sur les marchés secondaires, et que l'émission comme le rachat étaient contraints
par les horaires bancaires.
\end{dated}

\section{Tutoriel~: les frais de base et le coût d'une transaction}\label{tut:m3:blockchains-for-traders:gas}

\begin{tutorial}
\textbf{Objectif.} Implémenter la règle des \omterm{def:m3:blockchains-for-traders:gas}{frais de base} dans l'arithmétique entière du protocole, la projeter sur des blocs pleins
et un choc de demande, et chiffrer le coût d'une transaction. \textbf{État final~:}
la \cref{fig:m3:blockchains-for-traders:shock} et les chiffres du problème du week-end.
\begin{tutsteps}
\item \textbf{La règle.} \Cref{prop:m3:blockchains-for-traders:basefee} telle que l'EIP l'écrit.
\omcode{code/firm/gasfee/firm_gasfee.py}{12}{35}{Les frais de base de l'EIP-1559, leur projection et le coût d'une transaction.}\label{lst:m3:blockchains-for-traders:basefee}
\item \textbf{Le choc.} La demande de \omterm{def:m3:blockchains-for-traders:gas}{gas} baisse avec les \omterm{def:m3:blockchains-for-traders:gas}{frais de base}~; les frais montent jusqu'à ce que la demande rejoigne la cible.
\omcode{code/markets-3/14-blockchains-for-traders/python/m3_chain.py}{20}{33}{Un choc de demande sous la règle des frais de base.}\label{lst:m3:blockchains-for-traders:shock}
\item \textbf{Exécuter} \texttt{m3\_chain.full\_blocks(20)}, \texttt{arbitrage\_breakeven()} et \texttt{fig\_chain.py}.
\end{tutsteps}
\textbf{Pour aller plus loin.} Rendre la demande du choc plus élastique et observer jusqu'où montent les \omterm{def:m3:blockchains-for-traders:gas}{frais de base}~; ajouter une
distribution des pourboires des concurrents et calculer le pourboire qui maximise le profit espéré d'un arbitrage.
\end{tutorial}

\section{Construction~: frais et délais de confirmation}\label{bld:m3:blockchains-for-traders:gasfee}

\begin{build}
\bfield{Rôle} Chaque action on-chain de la mini-entreprise (un transfert entre plateformes, un swap, une liquidation)
doit être chiffrée et datée avant d'être envoyée~: le composant de frais projette les \omterm{def:m3:blockchains-for-traders:gas}{frais de base} et tarifie les transactions~; le
modèle de confirmation indique au planificateur de transferts du \cref{ch:m3:spot-markets} combien de temps les fonds sont en transit.
\bfield{Interface} \texttt{next\_base\_fee(parent\_base\_fee, gas\_used, gas\_limit)} en wei~;
\texttt{project(base\_fee, usage, gas\_limit)}~; \texttt{tx\_cost\_usd(gas, base\_fee\_gwei, priority\_gwei,
eth\_usd)}~; \texttt{confirmation\_seconds(chain, confirmations)}.
\bfield{Règles} Arithmétique entière en wei, exactement comme dans l'EIP (division arrondie par défaut, hausse minimale d'un wei)~;
les chaînes inconnues provoquent une erreur, jamais une valeur par défaut.
\bfield{Tests d'acceptation} \texttt{code/firm/gasfee/tests/}~: $\pm12.5\%$ à blocs pleins et vides~; inchangé
à la cible~; le minimum d'un wei~; vingt blocs pleins donnent $1.125^{20}$~; délais de confirmation.
\bfield{Pour aller plus loin} Frais de blobs pour les données de \omterm{def:m3:blockchains-for-traders:rollup}{rollup}~; marchés des frais d'autres chaînes~; probabilité d'inclusion en fonction du
pourboire, à partir des \omterm{def:m3:blockchains-for-traders:mempool}{mempools} observés.
\end{build}

\begin{omsources}
\item S.~Nakamoto, \emph{Bitcoin: A Peer-to-Peer Electronic Cash System}, 2008.
\item V.~Buterin et al., \emph{EIP-1559: Fee market change for ETH 1.0 chain}~; ethereum.org, proof-of-stake and The
Merge documentation.
\item Bloc Ethereum 26\,048\,449 via un nœud public JSON-RPC~; L2BEAT, scaling summary, 24 septembre 2026.
\item Federal Reserve, FEDS Notes, «~In the shadow of bank runs~: lessons from the Silicon Valley Bank failure and its
impact on stablecoins'', 17 décembre 2025.
\item CoinDesk, ``US officials tie North Korea's Lazarus hackers to \$625M Axie Infinity Ronin exploit'', 14 avril
2022.
\item Offre de stablecoins~: StablecoinBeat tracker, septembre 2026.
\end{omsources}

\section{Exercices}

\begin{exercise}[$\star$]\label{exo:m3:blockchains-for-traders:1}
Les \omterm{def:m3:blockchains-for-traders:gas}{frais de base} sont de 20~gwei et le bloc parent a utilisé 75~\% de sa limite. Donner les \omterm{def:m3:blockchains-for-traders:gas}{frais de base} du bloc suivant.
\end{exercise}

\begin{exercise}[$\star$]\label{exo:m3:blockchains-for-traders:2}
Combien de temps une plateforme qui exige six confirmations Bitcoin fait-elle attendre un déposant, en moyenne~? Et pour un
dépôt Ethereum crédité à la \omterm{def:m3:blockchains-for-traders:finality}{finalité}~?
\end{exercise}

\begin{exercise}[$\star$]\label{exo:m3:blockchains-for-traders:3}
Pourquoi l'ancrage d'un \omterm{def:m3:blockchains-for-traders:stable}{stablecoin} dépend-il de qui peut le racheter~?
\end{exercise}

\begin{exercise}[$\star\star$]\label{exo:m3:blockchains-for-traders:4}
Combien de blocs pleins consécutifs doublent les \omterm{def:m3:blockchains-for-traders:gas}{frais de base}~? Combien de blocs vides les divisent par deux~?
\end{exercise}

\begin{exercise}[$\star\star$]\label{exo:m3:blockchains-for-traders:5}
Un swap utilise 150\,000 \omterm{def:m3:blockchains-for-traders:gas}{gas}~; les \omterm{def:m3:blockchains-for-traders:gas}{frais de base} sont de 30~gwei, le pourboire de 2~gwei et l'ether vaut 3\,000~\$. Combien coûte-t-il~?
\end{exercise}

\begin{exercise}[$\star\star$]\label{exo:m3:blockchains-for-traders:6}
Expliquer pourquoi une transaction sur un \omterm{def:m3:blockchains-for-traders:rollup}{rollup} peut être bon marché et rapide sans être finale.
\end{exercise}

\begin{exercise}[$\star\star\star$]\label{exo:m3:blockchains-for-traders:7}
\emph{Programmation.} Avec \texttt{demand\_shock}, trouver le niveau le plus haut atteint par les \omterm{def:m3:blockchains-for-traders:gas}{frais de base} et le bloc à partir duquel
les blocs cessent pour la première fois d'être pleins.
\end{exercise}

\begin{exercise}[$\star\star\star$]\label{exo:m3:blockchains-for-traders:8}
\emph{Trouver la faille.} «~Un \omterm{def:m3:blockchains-for-traders:stable}{stablecoin} adossé un pour un à des bons du Trésor ne peut pas perdre son ancrage.~»
\end{exercise}

\section{Problème~: le pic de frais}

\begin{problem}[{Problème du week-end --- un arbitrage en concurrence sur une chaîne congestionnée}]\label{pb:m3:blockchains-for-traders:1}
Un searcher exécute un arbitrage qui rapporte 60~\$ avant \omterm{def:m3:blockchains-for-traders:gas}{gas}, utilise 150\,000 \omterm{def:m3:blockchains-for-traders:gas}{gas} et verse un pourboire de 2~gwei. Les \omterm{def:m3:blockchains-for-traders:gas}{frais de base} sont de 10~gwei et
l'ether coûte 3\,000~\$ (à titre illustratif). Un lancement de \omterm{def:m3:blockchains-for-traders:key}{jeton} très suivi remplit chaque bloc.

\medskip\noindent\textbf{Partie I --- Les frais.}
\begin{enumerate}
\item Combien l'arbitrage paie-t-il en \omterm{def:m3:blockchains-for-traders:gas}{gas} maintenant~?
\item Quels sont les \omterm{def:m3:blockchains-for-traders:gas}{frais de base} après un bloc plein, et après vingt~?
\item Combien l'arbitrage paie-t-il ensuite en \omterm{def:m3:blockchains-for-traders:gas}{gas}~?
\item Combien de temps prennent vingt blocs~?
\item Qui reçoit les \omterm{def:m3:blockchains-for-traders:gas}{frais de base}, et qui reçoit le pourboire~?
\end{enumerate}

\medskip\noindent\textbf{Partie II --- Le point mort.}
\begin{enumerate}[resume]
\item À quel \omterm{def:m3:blockchains-for-traders:gas}{frais de base} l'arbitrage cesse-t-il d'être rentable~?
\item Au bout de combien de blocs pleins ce seuil est-il atteint~?
\item Que se passe-t-il si l'arbitrage rapporte 120~\$~?
\item Pourquoi le searcher pourrait-il relever son pourboire pendant le pic~?
\item Que devient le \omterm{def:m3:blockchains-for-traders:gas}{frais de base} lorsque le lancement prend fin~?
\end{enumerate}

\medskip\noindent\textbf{Partie III --- Solutions de rechange.}
\begin{enumerate}[resume]
\item Un \omterm{def:m3:blockchains-for-traders:rollup}{rollup} serait-il moins cher, et à quoi le searcher renoncerait-il~?
\item Comment une soumission privée modifie-t-elle la concurrence entre searchers~?
\item Pourquoi les \omterm{def:m3:blockchains-for-traders:gas}{frais de base} sont-ils brûlés plutôt que versés au proposant~?
\item Qu'est-ce que le pic change pour les utilisateurs qui ne sont pas arbitragistes~?
\item Qu'apprend le searcher du \omterm{def:m3:blockchains-for-traders:mempool}{mempool} pendant le pic~?
\end{enumerate}

\medskip\noindent\textbf{Partie IV --- Jugement.}
\begin{enumerate}[resume]
\item Le marché des frais est-il un bon moyen d'allouer l'espace des blocs~?
\item De quoi un trader qui transfère des fonds entre plateformes a-t-il besoin dans ce chapitre~?
\item Pourquoi les frais comptent-ils davantage pour les petits arbitrages que pour les grands~?
\item Énoncer le \emph{résultat nommé}~: le multiplicateur des \omterm{def:m3:blockchains-for-traders:gas}{frais de base} après vingt blocs pleins, et le nombre de blocs pleins
au bout duquel l'arbitrage cesse d'être rentable.
\item En une phrase~: qu'achètent les frais d'une transaction~?
\end{enumerate}
\end{problem}

\section{Questions d'entretien}

\begin{interviewq}[$\star$ \iqroles{trader}]\label{iq:m3:blockchains-for-traders:1}
Qu'est-ce que la \omterm{def:m3:blockchains-for-traders:finality}{finalité}, et pourquoi compte-t-elle pour transférer des fonds entre bourses~?
\end{interviewq}

\begin{interviewq}[$\star$ \iqroles{développeur}]\label{iq:m3:blockchains-for-traders:2}
Expliquez le \omterm{def:m3:blockchains-for-traders:gas}{frais de base} et le \omterm{def:m3:blockchains-for-traders:gas}{frais de priorité} de l'EIP-1559.
\end{interviewq}

\begin{interviewq}[$\star\star$ \iqroles{chercheur, trader}]\label{iq:m3:blockchains-for-traders:3}
Comment un \omterm{def:m3:blockchains-for-traders:stable}{stablecoin} adossé à des réserves maintient-il son ancrage, et comment peut-il échouer~?
\end{interviewq}

\begin{interviewq}[$\star\star$ \iqroles{risque}]\label{iq:m3:blockchains-for-traders:4}
Quels sont les risques liés à la détention d'actifs transférés d'une chaîne à une autre~?
\end{interviewq}

\begin{interviewq}[$\star\star$ \iqroles{développeur}]\label{iq:m3:blockchains-for-traders:5}
Un dépôt n'affiche qu'une confirmation. Quand votre système doit-il le créditer, et que peut-il mal tourner~?
\end{interviewq}

\begin{interviewq}[$\star\star\star$ \iqroles{développeur, chercheur}]\label{iq:m3:blockchains-for-traders:6}
Concevez un estimateur de frais qui indique à un système de trading quel pourboire payer pour être inclus dans deux blocs avec
une probabilité de 95~\%.
\end{interviewq}
